\chapter{Distributed C++}

A distributed system is one where a message can be lost, delayed, duplicated, or reordered, and where any participant can fail independently while the others keep running. Moving C++ from a single process into this world changes the cost model by orders of magnitude --- a call is nanoseconds, a round trip is microseconds to milliseconds --- and introduces failure modes with no single-process analogue. This chapter covers the three pillars: serialization, remote procedure calls, and consensus.

\section*{Chapter Roadmap}
\begin{itemize}
\item 84.1 The Distributed Cost Model and Failure Model
\item 84.2 Serialization: Packing Data for the Wire
\item 84.3 Remote Procedure Calls
\item 84.4 RPC Delivery Semantics and Idempotency
\item 84.5 Consensus: Agreeing Under Failure (Raft)
\item 84.6 Hazards: Partial Failure, Split Brain, and the Fallacies
\end{itemize}

\section{The Distributed Cost Model and Failure Model}
A local call is $\sim$1\,ns; memory $\sim$100\,ns; a same-datacenter round trip $\sim$50--500\,$\mu$s; cross-continent $\sim$100+\,ms --- the network is 5--8 orders of magnitude slower than a call. The failure model adds a third outcome beyond return/crash: \textbf{partial failure} --- the call was sent but you never learn whether it arrived, executed, or replied.

\textbf{Why this matters.} The latency gap means batch, pipeline, and co-locate. Partial failure means every remote operation must define timeout, retry, and duplicate behaviour. Treating remote calls like local calls (the fallacy ``the network is reliable'') corrupts data the first time a packet drops.

\section{Serialization: Packing Data for the Wire}
\begin{lstlisting}[language=C++, caption={A length-prefixed binary serializer. Min standard: C++17.}]
class Buffer {
    std::vector<uint8_t> data_;
public:
    template <typename T>
    void write(const T& val) {
        static_assert(std::is_trivially_copyable_v<T>);
        const auto* p = reinterpret_cast<const uint8_t*>(&val);
        data_.insert(data_.end(), p, p + sizeof(T));
    }
    void write_string(const std::string& s) {
        write<uint32_t>(static_cast<uint32_t>(s.size()));
        const auto* p = reinterpret_cast<const uint8_t*>(s.data());
        data_.insert(data_.end(), p, p + s.size());
    }
    const uint8_t* begin() const { return data_.data(); }
    size_t size() const { return data_.size(); }
};
\end{lstlisting}

\textbf{Why this matters / cost model.} The \texttt{is\_trivially\_copyable} check gates the bit-copy: you may copy a POD but not a \texttt{std::string} (serialize length + bytes). Real systems handle endianness and schema evolution, hence Protobuf/FlatBuffers/Cap'n Proto/SBE. FlatBuffers/Cap'n Proto allow zero-copy access --- right for low-latency paths where parse cost dominates.

\section{Remote Procedure Calls}
\begin{lstlisting}[language=C++, caption={An RPC client stub; the function ID selects the handler. Min standard: C++11.}]
int Add(int a, int b) {
    Buffer buf;
    buf.write<uint32_t>(101);          // function ID for 'Add'
    buf.write(a); buf.write(b);
    return network.send_and_wait(buf); // serialize, send, block, deserialize
}
\end{lstlisting}

\textbf{Why this matters.} RPC structures a distributed system as ordinary calls, but the abstraction hides the network, tempting you to forget the latency and partial failure. A local \texttt{Add} cannot time out, retry, or run twice; the remote one can. gRPC/Thrift generate stubs from an IDL and surface errors. Treat every RPC as returning ``result or failure.''

\section{RPC Delivery Semantics and Idempotency}
\begin{itemize}
\item At-most-once --- never retry; may lose operations.
\item At-least-once --- retry until acked; may execute more than once.
\item Exactly-once --- appears to happen once; needs dedup + state.
\end{itemize}

\textbf{Why this matters.} If a client sends ``transfer \$100,'' times out, and retries, did the first execute? With at-least-once the retry may double it. The fix is idempotency: design so executing twice equals once (``set balance to X,'' or a deduplicated request ID). Exactly-once is constructed from at-least-once plus idempotency, not provided by the network. Every retry path must answer: safe to execute twice?

\section{Consensus: Agreeing Under Failure (Raft)}
\begin{lstlisting}[language=C++, caption={Raft leader-election state. Min standard: C++11.}]
enum class State { Follower, Candidate, Leader };
struct Node {
    State state = State::Follower;
    int current_term = 0; int voted_for = -1; int my_id;
    void on_election_timeout() {
        if (state == State::Follower) {
            state = State::Candidate;
            ++current_term; voted_for = my_id;
            request_votes();
        }
    }
    void request_votes();
};
\end{lstlisting}

\textbf{Why this matters / cost model.} Consensus gives a single source of truth surviving a minority of failures. Committing requires a majority round trip; a 5-node cluster tolerates 2 failures. The term number prevents two leaders acting on stale info --- a higher term wins, so a partitioned old leader is superseded. Raft/Paxos underpin etcd, ZooKeeper, CockroachDB. Use consensus for small critical metadata, not the bulk data path.

\section{Hazards: Partial Failure, Split Brain, and the Fallacies}
The fallacies: the network is reliable; latency is zero; bandwidth is infinite; the network is secure; topology doesn't change. Two dominant hazards: partial failure (a node may have executed your request despite your timeout --- design for ``don't know'') and split brain (a partition leaves two halves each believing it leads; quorum consensus prevents divergence).

\textbf{The discipline.} Distributed C++ is single-process C++ plus an adversarial network and independently-failing components. The volume's performance disciplines still apply beneath a failure-aware layer: idempotent operations, explicit timeouts, quorum consensus for truth, and the assumption that anything that can fail will.
